\chapter{总结}
\label{chap:conclusion}

\sysname 用独立的 rate/window gate 表示预装 subchannel，并以反馈门控选择器逐
packet offset 分配路径。用于在网 fan-in 时，组播 credit 把所选 subchannel 和
逐级聚合位置绑定到整组 requester，使通信组内的输入共享 Credit Context。
responder 根据整组完成事件控制后续 credit，并在超时或 owner mismatch 时启动
整组 repair。该传输对象使 packet-level 路径选择、聚合状态、拥塞反馈和恢复
使用相同边界。
六种集合通信原语只改变 channel 集合和 payload 方向，传输与恢复规则保持不变。

Normal credit 在带内推进各级共享循环 allocator，并把位置压栈返回；加速器
从 request 的位置栈直接访问共享 slot，并用 owner 隔离轮转后的迟到输入。
Normal request 单注入使数据面只需计数；requester 身份和 repair 去重保留在
responder。循环 allocator 的根本条件是 $A_e\geq D_e$，其中 $D_e$ 包含同一
domain 的全部 normal advance。严格 all-advance envelope 给出
$\lceil B_e+\lambda_eR_e\rceil$ 充分界；
把每个 fixed-$P$ data context 与一份 payload packet 一一对应，可得到通用
$2R_e$ 端口界，full-payload carrier 再收紧为 $R_e$。端口界对实际推进同一 allocator 的端口
求和，不随 job 数或 rank 数额外放大。
当 $A_e=D_e$ 时，带内循环分配以最小安全深度运行，并使 100\% 的已配置地址
成为无提前复用冲突的有效容量，不受哈希冲突或 per-job 分区影响。

包级仿真验证了共享状态、反馈驱动的逐-offset 选树、四种端侧拥塞控制算法和
故障恢复。在四棵候选树依次受限的
单-job 实验中，dynamic 达到 190.0~Gbps，而全程最优固定配额为
154.6~Gbps；前者在 4--5.5 个反馈周期内完成迁移。状态实验
分别覆盖共用端口和端口分散的 switch-wide allocator；110 次运行均满足相应
的循环复用界。在固定 32$\times$200-Gbps spine 和按端口界配置的 42-MiB
状态上，我们将配置人口从 16 增至 512，并令约 75\% 的 job 持续通信；相对
各点由拓扑确定的 per-job service budget，\sysname 保持 97.42\%--97.98\%，
静态分区的 SwitchML-like 在分区成为瓶颈后为 69.99\%--75.03\%，ATP-like
随哈希冲突增加从 97.59\% 降至 40.50\%。在 200-Gbps 一、二、三层拓扑
中，单-job 256-MiB AllReduce 分别达到 196.25、197.17 和
195.83~Gbps；50\%-active 的八-job 轨迹中，两/三层 aggregate goodput
为 763.96 和 754.97~Gbps，相对最优静态 SwitchML-like 提高
17.0\% 和 15.6\%。
当前 ns-3 模型不执行 FP32 运算，
1024-host 结果也只验证协议和状态路径。

% TODO(testbed-v80-final): 补最终四 rank V80 FP32、端侧 core、ATP/SwitchML、
% 资源/功耗与 GPU/PyTorch 训练片段结论；不要恢复旧三 rank DPDK proxy 数字。
